%% Official Springer Nature / Nature Portfolio submission template (sn-jnl).
%% Reference style: sn-nature (Nature Portfolio journals).
%% Build:  pdflatex main_sn && bibtex main_sn && pdflatex main_sn && pdflatex main_sn
\documentclass[sn-nature,lineno]{sn-jnl}

\usepackage{graphicx}
\usepackage{multirow}
\usepackage{amsmath,amssymb,amsfonts}
\usepackage{amsthm}
\usepackage{mathrsfs}
\usepackage[title]{appendix}
\usepackage{xcolor}
\usepackage{textcomp}
\usepackage{manyfoot}
\usepackage{booktabs}
\usepackage{longtable}
\usepackage{array}
\usepackage{calc}
\usepackage{algorithm}
\usepackage{algorithmicx}
\usepackage{algpseudocode}
\usepackage{listings}

% --- pdflatex compatibility for the pandoc-generated body ---
% (Unicode is substituted to LaTeX directly in body_sn.tex, so this
%  source compiles on any TeX Live / Overleaf with no extra packages.)
\providecommand{\tightlist}{\setlength{\itemsep}{0pt}\setlength{\parskip}{0pt}}
\providecommand{\pandocbounded}[1]{#1}

\theoremstyle{thmstyleone}
\newtheorem{theorem}{Theorem}
\newtheorem{proposition}[theorem]{Proposition}
\theoremstyle{thmstyletwo}
\newtheorem{example}{Example}
\newtheorem{remark}{Remark}
\theoremstyle{thmstylethree}
\newtheorem{definition}{Definition}

\raggedbottom

% Graphics path so \includegraphics{figXX.png} resolves
\graphicspath{{../outputs/figures/}}

\begin{document}

\title[Cross-disaster mapping is a calibration problem]{Cross-disaster mapping is largely a calibration problem, not a representation problem: four labels per event recover near-oracle accuracy and cut adaptation from days to minutes}

\author*[1]{\fnm{Qiming} \sur{Bao}}\email{qiming.bao@example.edu}
\author[2]{\fnm{Yanbing} \sur{Bai}}\email{ybbai@example.edu}

\affil*[1]{\orgdiv{[Department TBD]}, \orgname{[Institution TBD]},
  \orgaddress{\city{[City]}, \country{[Country]}}}
\affil[2]{\orgdiv{[Department TBD]}, \orgname{[Institution TBD]},
  \orgaddress{\city{[City]}, \country{[Country]}}}

\abstract{Deep-learning disaster mappers degrade sharply on events they were not trained on, and the field has treated this as a representation problem --- larger backbones, foundation models, multi-modal fusion. Across three public benchmarks spanning floods, building damage and wildfires (18 real events), we show the dominant mechanism is instead \textbf{calibration drift}: the pixel ranking transfers, but the decision threshold does not (15 of 16 measured event-optimal thresholds differ from the 0.5 default). Re-fitting one threshold recovers up to +0.235 F1 per event; none of three frozen foundation models (AlphaEarth, Prithvi, DOFA) exceeds a from-scratch U-Net. Four labelled tiles recover 99 \% of the full-pool oracle regardless of how they are chosen, whereas zero-label corrections (EM, BBSE) fail --- the class-conditional scores themselves distort, which only labels reveal. Because perception stays frozen, each new event costs four labels and minutes, not tens and days --- an order-of-magnitude cut in the only human step that scales with events, which an end-to-end agent turns into responder briefings. \emph{(150 words)}}

\keywords{cross-disaster generalisation, calibration drift, label shift,
foundation models, active learning, flood mapping, disaster response}

\maketitle

\input{body_sn}

\backmatter

\bmhead{Supplementary information}
The within-event-protocol ablations (sample-efficiency sweep;
decision-aligned reward A/B) and the full per-fold result JSONs are
provided as Supplementary Information.

\bmhead{Acknowledgements}
[TBD]

\section*{Declarations}
\begin{itemize}
\item \textbf{Funding}: [TBD]
\item \textbf{Competing interests}: The authors declare no competing
  interests.
\item \textbf{Data availability}: All benchmark data are public
  (Sen1Floods11, xBD, HLS Burn-Scars). All intermediate results,
  figures, and code are available at
  \url{https://github.com/14H034160212/geodisaster-fm}.
\item \textbf{Code availability}: As above; a versioned release with a
  Zenodo DOI will accompany publication.
\item \textbf{Author contributions}: [TBD]
\end{itemize}

\clearpage
\input{figures_sn}

\bibliography{../refs}

\end{document}
